
\documentclass[aps,preprint]{revtex4}
%%%%%%%%%%%%%%%%%%%%%%%%%%%%%%%%%%%%%%%%%%%%%%%%%%%%%%%%%%%%%%%%%%%%%%%%%%%%%%%%%%%%%%%%%%%%%%%%%%%%%%%%%%%%%%%%%%%%%%%%%%%%
\usepackage{amssymb}
\usepackage{amsmath}

\setcounter{MaxMatrixCols}{10}
%TCIDATA{OutputFilter=LATEX.DLL}
%TCIDATA{Version=4.10.0.2363}
%TCIDATA{Created=Sunday, February 03, 2002 11:33:45}
%TCIDATA{LastRevised=Thursday, July 29, 2004 15:18:59}
%TCIDATA{<META NAME="GraphicsSave" CONTENT="32">}
%TCIDATA{<META NAME="DocumentShell" CONTENT="Articles\SW\REVTeX 4">}
%TCIDATA{CSTFile=revtex4.cst}

\newtheorem{theorem}{Theorem}
\newtheorem{acknowledgement}[theorem]{Acknowledgement}
\newtheorem{algorithm}[theorem]{Algorithm}
\newtheorem{axiom}[theorem]{Axiom}
\newtheorem{claim}[theorem]{Claim}
\newtheorem{conclusion}[theorem]{Conclusion}
\newtheorem{condition}[theorem]{Condition}
\newtheorem{conjecture}[theorem]{Conjecture}
\newtheorem{corollary}[theorem]{Corollary}
\newtheorem{criterion}[theorem]{Criterion}
\newtheorem{definition}[theorem]{Definition}
\newtheorem{example}[theorem]{Example}
\newtheorem{exercise}[theorem]{Exercise}
\newtheorem{lemma}[theorem]{Lemma}
\newtheorem{notation}[theorem]{Notation}
\newtheorem{problem}[theorem]{Problem}
\newtheorem{proposition}[theorem]{Proposition}
\newtheorem{remark}[theorem]{Remark}
\newtheorem{solution}[theorem]{Solution}
\newtheorem{summary}[theorem]{Summary}
\newenvironment{proof}[1][Proof]{\noindent\textbf{#1.} }{\ \rule{0.5em}{0.5em}}
\input{tcilatex}

\begin{document}

\preprint{HEP/123-qed}
\title[Short title for running header]{Long title that appears on the title
page}
\author{First Author}
\affiliation{My Institution}
\author{Second Author}
\affiliation{My Institution}
\author{Third Author}
\affiliation{Other Institution}
\keywords{one two three}
\pacs{PACS number}

\begin{abstract}
Shell document for REV\TeX{} 4.
\end{abstract}

\volumeyear{year}
\volumenumber{number}
\issuenumber{number}
\eid{identifier}
\date[Date text]{date}
\received[Received text]{date}
\revised[Revised text]{date}
\accepted[Accepted text]{date}
\published[Published text]{date}
\startpage{101}
\endpage{102}
\maketitle
\tableofcontents

If one considers subgroups $\left\{ \left. SU\left( N_{1}\right) ,\cdots
,SU\left( N_{m}\right) \right\vert \tsum_{1\leq j\leq m}N_{j}=N\right\} $ of 
$SU\left( N\right) $ and $\left\{ \left. SU\left( \nu _{1}\right) ,\cdots
,SU\left( \nu _{m^{\prime }}\right) \right\vert \tsum_{1\leq j\leq m^{\prime
}}\nu _{j}=2s-N\right\} $%
\begin{equation}
\mathcal{A}\left( \Psi ,\tau \right) =\int_{t_{i}}^{t_{f}}\mathcal{L}\left(
\Psi ,\dot{\Psi},t\right) dt
\end{equation}%
\begin{equation}
\delta \mathcal{A}\left( \Psi ,\tau \right) =\int_{t_{i}}^{t_{f}}\delta 
\mathcal{L}\left( \Psi ,\dot{\Psi},t\right) dt
\end{equation}%
\begin{eqnarray}
\delta \mathcal{L}\left( \Psi ,\dot{\Psi},t\right) &=&\int \left\{ \frac{%
\delta \mathcal{L}\left( \Psi ,\dot{\Psi},t\right) }{\delta \Psi \left( 
\mathsf{z}\right) }\delta \Psi \left( \mathsf{z}\right) +\frac{\delta 
\mathcal{L}\left( \Psi ,\dot{\Psi},t\right) }{\delta \dot{\Psi}\left( 
\mathsf{z}\right) }\delta \dot{\Psi}\left( \mathsf{z}\right) \right. \\
&&\left. +\frac{\delta \mathcal{L}\left( \Psi ,\dot{\Psi},t\right) }{\delta
\Psi ^{\ast }\left( \mathsf{z}\right) }\delta \Psi ^{\ast }\left( \mathsf{z}%
\right) +\frac{\delta \mathcal{L}\left( \Psi ,\dot{\Psi},t\right) }{\delta 
\dot{\Psi}^{\ast }\left( \mathsf{z}\right) }\delta \dot{\Psi}^{\ast }\left( 
\mathsf{z}\right) \right\} d\mathsf{z}
\end{eqnarray}%
\begin{eqnarray}
\mathcal{L}\left( \Psi ,\dot{\Psi},t\right) &\equiv &\mathcal{L}\left( \Psi
\left( t\right) \right) \\
&\equiv &\mathcal{L}\left( \Psi \left( t\right) ,\dot{\Psi}\left( t\right)
\right) \equiv \mathcal{L}\left( \Psi ,\dot{\Psi}\right)
\end{eqnarray}%
\begin{equation}
\mathcal{L}\left( \Psi ,\dot{\Psi},t\right) =\frac{\left\langle \Psi \left(
t\right) |\left( i\frac{\partial }{\partial t}-H\right) \Psi \left( t\right)
\right\rangle }{\left\langle \Psi \left( t\right) |\Psi \left( t\right)
\right\rangle }
\end{equation}

\begin{eqnarray}
\delta \mathcal{L}\left( \Psi \left( t\right) \right) &=&\frac{1}{%
\left\langle \Psi \left( t\right) |\Psi \left( t\right) \right\rangle }%
\left\{ \left\langle \delta \Psi \left( t\right) |\left( i\frac{\partial }{%
\partial t}-H\right) \Psi \left( t\right) \right\rangle +\left\langle \Psi
\left( t\right) |\left( i\frac{\partial }{\partial t}-H\right) \delta \Psi
\left( t\right) \right\rangle \right.  \notag \\
&&\qquad \qquad \qquad \qquad \qquad \qquad \left. -\mathcal{L}\left( \Psi
\left( t\right) \right) \left\{ \left\langle \delta \Psi \left( t\right)
|\Psi \left( t\right) \right\rangle +\left\langle \Psi \left( t\right)
|\delta \Psi \left( t\right) \right\rangle \right\} \right\}
\end{eqnarray}

\begin{equation}
\tau =\left[ t_{i},t_{f}\right]
\end{equation}

\begin{equation}
\delta \Psi \left( t_{f}\right) =\delta \Psi \left( t_{i}\right) =0
\end{equation}

\begin{equation}
\delta \Psi \left( t_{f}\right) =\delta \Psi \left( t_{i}\right) =0
\end{equation}

\begin{equation}
\tau =\left[ t_{i},t_{f}\right]
\end{equation}

\begin{equation}
\frac{\delta \mathcal{L}\left( \Psi \left( t\right) \right) }{\delta \dot{%
\Psi}^{\ast }\left( \mathsf{z}\right) }=0
\end{equation}%
\begin{equation}
\delta _{\Psi }\mathcal{A}\left( \Psi ,\tau \right) =0
\end{equation}%
\begin{equation}
\left( i\frac{\partial }{\partial t}-H\right) \left| \Psi \left( t\right)
\right\rangle =\left( \mathcal{L}\left( \Psi \left( t\right) \right) -i\frac{%
d}{dt}\ln \left\langle \Psi \left( t\right) |\Psi \left( t\right)
\right\rangle \right) ^{\ast }\left| \Psi \left( t\right) \right\rangle
\end{equation}%
\begin{equation}
\left( i\frac{\partial }{\partial t}-H\right) \left| \Psi \left( t\right)
\right\rangle =\mathcal{L}\left( \Psi \left( t\right) \right) \left| \Psi
\left( t\right) \right\rangle
\end{equation}%
\begin{eqnarray}
\left\langle \left[ i\frac{\partial }{\partial t}\right] ^{\dagger }\varphi
_{1}\left( t\right) |\varphi _{2}\left( t\right) \right\rangle &=&i\frac{d}{%
dt}\left\langle \varphi _{1}\left( t\right) |\varphi _{2}\left( t\right)
\right\rangle +\left\langle i\frac{\partial }{\partial t}\varphi _{1}\left(
t\right) |\varphi _{2}\left( t\right) \right\rangle \\
&&\qquad \qquad \qquad \qquad \qquad \qquad \quad \forall \varphi
_{1},\varphi _{2},t
\end{eqnarray}%
\begin{equation}
i\delta \Psi \left( t\right)
\end{equation}%
\begin{equation}
\delta \Psi \left( t\right) \qquad \qquad \delta \Psi ^{\ast }\left( t\right)
\end{equation}%
\begin{equation}
\Phi \left( \mathsf{z,}t\right) =\Psi \left( \mathsf{z,}t\right)
e^{i\int_{t_{i}}^{t}\mathcal{L}\left( \Psi ,\dot{\Psi},t^{\prime }\right)
dt^{\prime }}=\Psi \left( \mathsf{z,}t\right) e^{i\mathcal{A}\left( \Psi ,%
\left[ t_{i},t\right] \right) }
\end{equation}%
\begin{equation}
\dot{\Phi}\left( \mathsf{z,}t\right) =\dot{\Psi}\left( \mathsf{z,}t\right)
e^{i\int_{t_{i}}^{t}\mathcal{L}\left( \Psi ,\dot{\Psi},t^{\prime }\right)
dt^{\prime }}+i\Psi \left( \mathsf{z,}t\right) \mathcal{L}\left( \Psi ,\dot{%
\Psi},t\right) e^{i\int_{t_{i}}^{t}\mathcal{L}\left( \Psi ,\dot{\Psi}%
,t^{\prime }\right) dt^{\prime }}
\end{equation}%
\begin{eqnarray}
&=&\dot{\Psi}\left( \mathsf{z,}t\right) e^{i\int_{t_{i}}^{t}\mathcal{L}%
\left( \Psi ,\dot{\Psi},t^{\prime }\right) dt^{\prime }}  \notag \\
&&+i\left\{ i\frac{\partial \Psi \left( \mathsf{z,}t\right) }{\partial t}%
-\int H\left( \mathsf{z,z}^{\prime }\right) \Psi \left( \mathsf{z}^{\prime }%
\mathsf{,}t\right) d\mathsf{z}^{\prime }\right\} e^{i\int_{t_{i}}^{t}%
\mathcal{L}\left( \Psi ,\dot{\Psi},t^{\prime }\right) dt^{\prime }}  \notag
\end{eqnarray}%
\begin{equation}
=-i\left[ H\Psi \right] \left( \mathsf{z,}t\right) e^{-i\int_{t_{i}}^{t}%
\mathcal{L}\left( \Psi ,\dot{\Psi},t^{\prime }\right) dt^{\prime }}
\end{equation}%
\begin{equation}
\mathcal{L}\left( \Psi ,\dot{\Psi},t\right)
\end{equation}%
\begin{equation}
\mathcal{A}\left( \Psi ,\tau \right) =\int_{t_{i}}^{t_{f}}\mathcal{L}\left(
\Psi ,\dot{\Psi},t\right) dt
\end{equation}%
\begin{equation}
\mathcal{L}\left( \Psi ,\dot{\Psi},t\right) =\frac{\left\langle \Psi \left(
t\right) |\left( i\frac{\partial }{\partial t}-H\right) \Psi \left( t\right)
\right\rangle }{\left\langle \Psi \left( t\right) |\Psi \left( t\right)
\right\rangle }  \notag
\end{equation}%
\begin{eqnarray}
&\Rightarrow &\dot{\Phi}\left( \mathsf{z,}t\right) =-iH\Phi \left( \mathsf{z,%
}t\right) \equiv i\frac{\partial \Phi \left( \mathsf{z,}t\right) }{\partial t%
}=H\Phi \left( \mathsf{z,}t\right) \equiv \left( i\frac{\partial \Phi \left( 
\mathsf{z,}t\right) }{\partial t}-H\Phi \left( \mathsf{z,}t\right) \right) =0
\notag \\
&\equiv &\left\{ i\frac{\partial }{\partial t}-H\right\} \left| \Phi \left(
t\right) \right\rangle =0  \label{SE}
\end{eqnarray}%
\begin{equation}
\mathcal{L}\left( \Phi ,\dot{\Phi},t\right) =0\text{ }
\end{equation}%
\begin{equation}
\int \Phi ^{\ast }\left( \mathsf{z,}t\right) \Phi \left( \mathsf{z,}t\right)
d\mathsf{z}=\left\langle \Psi \left( t_{i}\right) |\Psi \left( t_{i}\right)
\right\rangle
\end{equation}%
\begin{equation}
\left| \Phi \left( t_{f}\right) \right\rangle \neq \left| \Psi \left(
t_{f}\right) \right\rangle
\end{equation}%
\begin{equation}
\left| \Phi \left( t_{i}\right) \right\rangle =\left| \Psi \left(
t_{i}\right) \right\rangle
\end{equation}%
\begin{equation}
\left. \equiv \right. \left| \Phi \left( t_{f}\right) \right\rangle =\left|
\Psi \left( t_{f}\right) \right\rangle e^{i\mathcal{A}\left( \Psi ,\tau
\right) }=\left( \frac{\left\langle \Psi \left( t_{i}\right) |\Psi \left(
t_{i}\right) \right\rangle }{\left\langle \Psi \left( t_{f}\right) |\Psi
\left( t_{f}\right) \right\rangle }\right) ^{\frac{1}{2}}\left| \Psi \left(
t_{f}\right) \right\rangle e^{i\func{Re}\mathcal{A}\left( \Psi ,\tau \right)
}
\end{equation}%
\begin{equation}
\Phi \left( \mathsf{z,}t_{f}\right) =\left( \frac{\left\langle \Psi \left(
t_{i}\right) |\Psi \left( t_{i}\right) \right\rangle }{\left\langle \Psi
\left( t_{f}\right) |\Psi \left( t_{f}\right) \right\rangle }\right) ^{\frac{%
1}{2}}\Psi \left( \mathsf{z,}t_{f}\right) e^{i\func{Re}\mathcal{A}\left(
\Psi ,\tau \right) }  \label{phi}
\end{equation}%
and in general 
\begin{equation}
\func{Re}\mathcal{A}\left( \Psi ,\tau \right) \neq 0
\end{equation}%
We can see from eq. $\left( \ref{phi}\right) $ that the norm of the final
state is fixed by the choice of initial ''state'' but the phase is not. The
equation that determines the final phase is 
\begin{equation}
\frac{d\sigma }{dt}=-i\left\{ \frac{\func{Im}\left\langle \Psi \left(
t\right) |\dot{\Psi}\left( t\right) \right\rangle }{\left\langle \Psi \left(
t\right) |\Psi \left( t\right) \right\rangle }+\frac{\left\langle \Psi
\left( t\right) |H\Psi \left( t\right) \right\rangle }{\left\langle \Psi
\left( t\right) |\Psi \left( t\right) \right\rangle }\right\} =\func{Re}%
\mathcal{L}\left( \Psi \left( t\right) ,\dot{\Psi}\left( t\right) \right)
;\sigma \left( t_{i}\right) =0  \label{phase}
\end{equation}%
where $\left| \Psi \left( t\right) \right\rangle $ is a solution of eq. $%
\left( \ref{preSE}\right) $ and the solution of eq. $\left( \ref{phase}%
\right) $ is 
\begin{equation}
\sigma \left( t\right) =\func{Re}\mathcal{A}\left( \Psi ,\left[ t_{i},t%
\right] \right)
\end{equation}%
\begin{equation}
\frac{1}{2}\frac{d\ln \left\langle \Psi |\Psi \right\rangle }{dt}=\frac{1}{2}%
\frac{1}{\left\langle \Psi |\Psi \right\rangle }\left\{ \left\langle \dot{%
\Psi}|\Psi \right\rangle +\left\langle \Psi |\dot{\Psi}\right\rangle \right\}
\end{equation}%
\begin{eqnarray}
&&\mathcal{L}\left( \Psi \left( t\right) ,\dot{\Psi}\mathcal{S}_{0}\left(
t_{f}\right) =\left\{ \left| \delta \Psi \left( t\right) \right\rangle
\right\} ;\quad \left| \delta \Psi \left( t\right) \right\rangle =0,\quad
\kappa =i,f;\quad t\in \left[ t_{i},t_{f}\right] \left( t\right) \right) =%
\frac{i\left\langle \Psi |\dot{\Psi}\rangle -\left\langle \Psi \right| H\Psi
\right\rangle }{\left\langle \Psi |\Psi \right\rangle } \\
&&\mathcal{L}\left( \Psi \left( t\right) ,\dot{\Psi}\left( t\right) \right)
^{\ast }=\frac{-i\left\langle \dot{\Psi}|\Psi \rangle -\left\langle \Psi
\right| H\Psi \right\rangle }{\left\langle \Psi |\Psi \right\rangle } \\
&&\mathcal{L}\left( \Psi \left( t\right) ,\dot{\Psi}\left( t\right) \right) -%
\mathcal{L}\left( \Psi \left( t\right) ,\dot{\Psi}\left( t\right) \right)
^{\ast }=\frac{i\left\{ \langle \Psi |\dot{\Psi}\rangle +\langle \dot{\Psi}%
|\Psi \rangle \right\} }{\left\langle \Psi |\Psi \right\rangle }=2i\func{Im}%
\mathcal{L}\left( \Psi \left( t\right) ,\dot{\Psi}\left( t\right) \right) \\
&&\frac{1}{2}\frac{d\ln \left\langle \Psi |\Psi \right\rangle }{dt}=\func{Im}%
\mathcal{L}\left( \Psi ,\dot{\Psi},t\right)
\end{eqnarray}%
\begin{equation}
\left| \eta \left( t\right) \Psi \left( t\right) \right\rangle
\rightsquigarrow D_{\Psi }^{N}\left( t\right) \qquad \forall \eta \left(
t\right)
\end{equation}%
\begin{equation}
\left| \Psi \left( \eta ,t_{\kappa }\right) \right\rangle =\eta \left(
t_{\kappa }\right) \left| \Lambda _{\kappa }\right\rangle ;\qquad \kappa
=i,f;\quad \eta \left( t_{\kappa }\right) \in \mathbb{C}
\end{equation}%
\begin{equation}
\delta \mathcal{A}\left( \Psi ,\tau \right) =\delta _{\Psi }^{{}}\mathcal{A}%
\left( \Psi ,\tau \right) +\delta _{\tau }\mathcal{A}\left( \Psi ,\tau
\right) =0
\end{equation}%
\begin{eqnarray}
\mathcal{S}_{1}\left( t_{f}\right) &=&\left\{ \left| \delta \Psi \left(
t\right) \right\rangle \right\} ;\quad \left| \delta \Psi \left( t_{\kappa
}\right) \right\rangle =\tilde{\eta}\left( t_{\kappa }\right) \left| \Lambda
_{\kappa }\right\rangle \\
\eta \left( t\right) &\in &\mathbb{C},\,\kappa =i,f;\quad t\in \left[
t_{i},t_{f}\right]
\end{eqnarray}%
\begin{equation}
\mathcal{S}_{0}\left( t_{f}\right) =\left\{ \left| \delta \Psi \left(
t\right) \right\rangle \right\} ;\quad \left| \delta \Psi \left( t\right)
\right\rangle =0,\quad \kappa =i,f;\quad t\in \left[ t_{i},t_{f}\right]
\end{equation}%
\begin{equation}
\text{such that }\eta \left( t_{\kappa }\right) =\eta _{\kappa };\quad
\kappa =i,f
\end{equation}%
For $t\in \left( t_{i},t_{f}\right) $ there are no constraints on the state
vectors

Consider scalar trajectories $\eta \left( t\right) \in \mathbb{C}$ that are
arbitrary

twice differentiable complex valued functions $\forall t\in \left[
t_{i},t_{f}\right] $%
\begin{equation}
\text{Define }\Psi \left( \eta ,t\right) =\eta \left( t\right) \Psi \left(
t\right) \qquad \forall t\in \left[ t_{i},t_{f}\right]
\end{equation}%
for $t\in \left( t_{i},t_{f}\right) $ this does nothing as any variation is
allowed

We can replace the state vector variables $\Psi \left( t\right) $ by $\eta
\left( t\right) \Psi \left( t\right) $ in the Action%
\begin{equation}
\text{for }t=t_{i},t_{f}\text{ this describes the constraints }
\end{equation}

\begin{equation}
\mathcal{S}_{1}\left( t_{f}\right) =\left\{ \left| \delta \Psi \left(
t\right) \right\rangle \right\} ;\quad \left| \delta \Psi \left( t_{\kappa
}\right) \right\rangle =\tilde{\eta}\left( t_{\kappa }\right) \left| \Lambda
_{\kappa }\right\rangle ,\,\tilde{\eta}\left( t\right) \in \mathbb{C}%
,\,\kappa =i,f;\quad t\in \left[ t_{i},t_{f}\right]
\end{equation}%
\begin{equation}
\mathcal{S}_{1}\left( t_{f}\right) \supset \mathcal{S}_{0}\left(
t_{f}\right) \text{ as }\tilde{\eta}\left( t_{\kappa }\right) =0\text{ is
always a possible}
\end{equation}%
\begin{equation}
\mathcal{L}\left( \Psi \left( \eta ,t\right) ,\dot{\Psi}\left( \eta
,t\right) \right) =\mathcal{L}\left( \Psi \left( t\right) ,\dot{\Psi}\left(
t\right) \right) +i\frac{d\ln \eta \left( t\right) }{dt}
\end{equation}

Noting that the change of action induced by complex scaling

of the initial and final point is a function of the relative scaling

of these two points one can regard the action for the general

problem as a function of the final position, $\eta \left( t_{f}\right)
\left| \Lambda _{f}\right\rangle $, the

time interval, $\tau $, a \emph{fixed} initial point, $\Lambda _{\kappa }$,
and of stationary

trajectories that are solutions of the fixed end point problem.

\bigskip

The general equations are obtained by first observing that

the first order variation $\delta _{\Psi }^{\left( 1\right) }\mathcal{A}$ of
the action wrt changes

in stationary fixed point end point trajectories is%
\begin{equation}
\frac{\delta \mathcal{L}}{\delta \Psi }=\frac{d}{dt}\frac{\delta \mathcal{L}%
}{\delta \dot{\Psi}}
\end{equation}

\bigskip

\bigskip

When we fix the end point of the motion $\left| \Lambda _{f}\right\rangle $

but change the arrival time $t_{f}$ to $t_{f}+\delta t_{f}$, we need

to change $\left| \Psi \left( t_{f}\right) \right\rangle $ to $\left| \Psi
\left( t_{f}\right) \right\rangle -\left| \dot{\Psi}\left( t_{f}\right)
\right\rangle \delta t$

so that it arrives at $\left| \Lambda _{f}\right\rangle $ at time $%
t_{f}+\delta t_{f}$ i.e.

\begin{equation}
\left| \eta \left( t\right) \Psi \left( t\right) \right\rangle
\end{equation}

which must be added to the direct change

induced by the change of length of process

Adding these two together we obtain

that the first order variation of the action

wrt time durations $\tau $ is

The total first order change in the

generalized stationary expression

is then given by%
\begin{equation}
\frac{\delta \mathcal{L}\left( \eta \Psi ,\overset{\cdot }{\eta \Psi }%
,t\right) }{\delta \Psi \left( \mathsf{z}\right) }-\frac{d}{dt}\left\{ \frac{%
\delta \mathcal{L}\left( \eta \Psi ,\overset{\cdot }{\eta \Psi },t\right) }{%
\delta \dot{\Psi}\left( \mathsf{z}\right) }\right\} =0
\end{equation}%
\begin{equation}
\frac{\delta \mathcal{L}\left( \Psi ,\overset{\cdot }{\Psi },t\right) }{%
\delta \Psi \left( \mathsf{z}\right) }-\frac{d}{dt}\left\{ \frac{\delta 
\mathcal{L}\left( \Psi ,\overset{\cdot }{\Psi },t\right) }{\delta \dot{\Psi}%
\left( \mathsf{z}\right) }\right\} =0
\end{equation}%
\begin{equation}
\int_{t_{i}}^{t_{f}}\left\{ \left\langle \left. \left[ \frac{\delta \mathcal{%
L}\left( \eta \Psi ,\overset{\cdot }{\eta \Psi },t\right) }{\delta \Psi }%
\right] ^{\ast }\right\vert \delta \Psi \left( t\right) \right\rangle
+\left\langle \left. \left[ \frac{\delta \mathcal{L}\left( \eta \Psi ,%
\overset{\cdot }{\eta \Psi },t\right) }{\delta \dot{\Psi}}\right] ^{\ast
}\right\vert \delta \dot{\Psi}\left( t\right) \right\rangle \right\} dt
\end{equation}%
\begin{eqnarray}
&&\left. =\right. \int_{t_{i}}^{t_{f}}\left\{ \left\langle \left. \left[ 
\frac{d}{dt}\left\{ \frac{\delta \mathcal{L}\left( \eta \Psi ,\overset{\cdot 
}{\eta \Psi },t\right) }{\delta \dot{\Psi}}\right\} \right] ^{\ast
}\right\vert \delta \Psi \left( t\right) \right\rangle \right.  \notag \\
&&\qquad \qquad \qquad \qquad \qquad \qquad \qquad \qquad \qquad +\left.
\left\langle \left. \left[ \frac{\delta \mathcal{L}\left( \eta \Psi ,\overset%
{\cdot }{\eta \Psi },t\right) }{\delta \dot{\Psi}}\right] ^{\ast
}\right\vert \delta \dot{\Psi}\left( t\right) \right\rangle dt\right\}
\end{eqnarray}%
as the trajectories $\left\vert \eta \Psi \left( t\right) \right\rangle $
satisfy the Euler-Lagrange equation%
\begin{eqnarray}
&&\int_{t_{i}}^{t_{f}}\left\langle \left. \left[ \frac{\delta \mathcal{L}%
\left( \eta \Psi ,\overset{\cdot }{\eta \Psi },t\right) }{\delta \dot{\Psi}}%
\right] ^{\ast }\right\vert \delta \dot{\Psi}\left( t\right) \right\rangle
dt=\int_{t_{i}}^{t_{f}}\left\langle \left. \left[ \frac{\delta \mathcal{L}%
\left( \eta \Psi ,\overset{\cdot }{\eta \Psi },t\right) }{\delta \dot{\Psi}}%
\right] ^{\ast }\right\vert \frac{d\delta \Psi \left( t\right) }{dt}%
\right\rangle dt  \notag \\
&=&\left\langle \left. \left[ \frac{\delta \mathcal{L}\left( \eta \Psi ,%
\overset{\cdot }{\eta \Psi },t_{f}\right) }{\delta \dot{\Psi}}\right] ^{\ast
}\right\vert \delta \Psi \left( t_{f}\right) \right\rangle
-\int_{t_{i}}^{t_{f}}\left\langle \left. \frac{d}{dt}\left\{ \frac{\delta 
\mathcal{L}\left( \eta \Psi ,\overset{\cdot }{\eta \Psi },t\right) }{\delta 
\dot{\Psi}}\right\} ^{\ast }\right\vert \delta \Psi \left( t\right)
\right\rangle dt
\end{eqnarray}%
\begin{equation}
\delta _{\Psi }^{\left( 1\right) }\mathcal{A}\left( \Psi ,\tau \right)
=\left\langle \left. \left[ \frac{\delta \mathcal{L}\left( \eta \Psi ,%
\overset{\cdot }{\eta \Psi },t\right) }{\delta \dot{\Psi}}\right] ^{\ast
}\right\vert \delta \Psi \left( t_{f}\right) \right\rangle
\end{equation}%
\begin{equation}
-\left\langle \left. \left[ \frac{\delta \mathcal{L}\left( \eta \Psi ,%
\overset{\cdot }{\eta \Psi },t_{f}\right) }{\delta \dot{\Psi}}\right] ^{\ast
}\right\vert \dot{\Psi}\left( t_{f}\right) \delta t\right\rangle
\end{equation}%
\begin{equation}
\mathcal{L}\left( \eta \Psi ,\overset{\cdot }{\eta \Psi },t_{f}\right)
\delta t
\end{equation}%
\begin{equation}
\delta _{\tau }^{\left( 1\right) }\mathcal{A}\left( \Psi ,\tau \right) =%
\mathcal{L}\left( \eta \Psi ,\overset{\cdot }{\eta \Psi },t_{f}\right)
\delta t-\left\langle \left. \left[ \frac{\delta \mathcal{L}\left( \eta \Psi
,\overset{\cdot }{\eta \Psi },t_{f}\right) }{\delta \dot{\Psi}}\right]
^{\ast }\right\vert \dot{\Psi}\left( t_{f}\right) \delta t\right\rangle
\end{equation}%
\begin{eqnarray}
\delta ^{\left( 1\right) }\mathcal{A}\left( \Psi ,\tau \right)
&=&\left\langle \left. \left[ \frac{\delta \mathcal{L}\left( \eta \Psi ,%
\overset{\cdot }{\eta \Psi },t_{f}\right) }{\delta \dot{\Psi}}\right] ^{\ast
}\right\vert \delta \Psi \left( t_{f}\right) \right\rangle +\mathcal{L}%
\left( \eta \Psi ,\overset{\cdot }{\eta \Psi },t_{f}\right) \delta t  \notag
\\
&&\qquad \qquad \qquad \qquad \qquad \qquad -\left\langle \left. \left[ 
\frac{\delta \mathcal{L}\left( \eta \Psi ,\overset{\cdot }{\eta \Psi }%
,t_{f}\right) }{\delta \dot{\Psi}}\right] ^{\ast }\right\vert \dot{\Psi}%
\left( t_{f}\right) \delta t\right\rangle  \notag \\
\delta ^{\left( 1\right) }\mathcal{A}\left( \Psi ,\tau \right)
&=&\left\langle \left. \left[ \frac{\delta \mathcal{L}\left( \eta \Psi ,%
\overset{\cdot }{\eta \Psi },t_{f}\right) }{\delta \dot{\Psi}}\right] ^{\ast
}\right\vert \delta \Psi \left( t_{f}\right) \right\rangle -\left\langle
\left. \left[ \frac{\delta \mathcal{L}\left( \eta \Psi ,\overset{\cdot }{%
\eta \Psi },t_{f}\right) }{\delta \dot{\Psi}}\right] ^{\ast }\right\vert 
\frac{d\Psi \left( t_{f}\right) }{dt}\delta t\right\rangle  \notag \\
&&\qquad \qquad \qquad \qquad \qquad \qquad \qquad \qquad \qquad \qquad
\qquad \qquad +\mathcal{L}\left( \eta \Psi ,\overset{\cdot }{\eta \Psi }%
,t_{f}\right) \delta t  \notag \\
&=&\left\langle \left. \left[ \frac{\delta \mathcal{L}\left( \eta \Psi ,%
\overset{\cdot }{\eta \Psi },t_{f}\right) }{\delta \dot{\Psi}}\right] ^{\ast
}\right\vert \delta \Psi \left( t_{f}\right) \right\rangle -\left\langle
\left. \left[ \frac{\delta \mathcal{L}\left( \eta \Psi ,\overset{\cdot }{%
\eta \Psi },t_{f}\right) }{\delta \dot{\Psi}}\right] ^{\ast }\right\vert
\delta \Psi \left( t_{f}\right) \right\rangle  \notag \\
&&\qquad \qquad \qquad \qquad \qquad \qquad \qquad \qquad \qquad \qquad
\qquad \qquad +\mathcal{L}\left( \eta \Psi ,\overset{\cdot }{\eta \Psi }%
,t_{f}\right) \delta t
\end{eqnarray}%
\begin{equation}
\delta ^{\left( 1\right) }\mathcal{A}\left( \Psi ,\tau \right) =\mathcal{L}%
\left( \eta \Psi ,\overset{\cdot }{\eta \Psi },t_{f}\right) \delta t
\end{equation}%
\begin{equation}
\delta ^{\left( 1\right) }\mathcal{A}\left( \Psi ,\tau \right) =\mathcal{L}%
\left( \eta \left( t_{f}\right) \left\vert \Lambda _{f}\right\rangle ,\left.
\left\vert \dot{\Psi}\left( t_{f}\right) \right\rangle \right\vert
_{\left\vert \Psi \left( t_{f}\right) \right\rangle =\eta \left(
t_{f}\right) \left\vert \Lambda _{f}\right\rangle }\right) +i\frac{d\ln \eta
\left( t_{\kappa }\right) }{dt}
\end{equation}%
\begin{equation}
\frame{$\mathcal{L}\left( \left\vert \Lambda _{f}\right\rangle ,\left.
\left\vert \dot{\Psi}\left( t_{f}\right) \right\rangle \right\vert
_{\left\vert \Psi \left( t_{f}\right) \right\rangle =\left\vert \Lambda
_{f}\right\rangle }\right) +i\frac{d\ln \eta \left( t_{\kappa }\right) }{dt}$%
=0}
\end{equation}%
\begin{equation}
\mathcal{L}\left( \eta \Psi ,\overset{\cdot }{\eta \Psi },t_{f}\right) =%
\mathcal{L}\left( \Psi ,\overset{\cdot }{\Psi },t_{f}\right) +i\frac{d\ln
\eta \left( t_{f}\right) }{dt}
\end{equation}%
\begin{equation}
i\frac{d\ln \eta \left( t_{f}\right) }{dt}=-\mathcal{L}\left( \Psi ,\overset{%
\cdot }{\Psi },t_{f}\right)
\end{equation}%
\begin{equation}
\eta \left( t\right) =e^{i\mathcal{A}\left( \Psi ,\left[ t_{i},t\right]
\right) }
\end{equation}%
\begin{equation}
\delta ^{\left( 1\right) }\mathcal{A}\left( \Psi ,\tau \right)
=0\Longrightarrow
\end{equation}%
\begin{eqnarray}
&&\text{ \ \ \ \ \ \ \ by considering } \\
&&1.\int_{t_{u}}^{t_{f}}\left\{ \left\langle \delta \Psi \left( t\right)
\left\vert \left( i\frac{\partial }{\partial t}-H\right) \Psi \left(
t\right) \right. \right\rangle \right. \\
&&\qquad \left. {}\right. \ \ \ \ \ \ \ \ \ \ \ \ \ \ \ \ \ \ \left. -%
\mathcal{L}\left( \Psi \left( t\right) \right) \left\langle \delta \Psi
\left( t\right) \left\vert \Psi \left( t\right) \right. \right\rangle
\right\} dt=0\quad \forall \delta \Psi \left( t\right) \\
\qquad &&\text{that satisfy the fixed end point boundary conditions and } \\
&&2.\text{ }\left\langle \Psi \left( t\right) \left\vert \left( i\frac{%
\partial }{\partial t}-H\right) \mathsf{z}\right. \right\rangle \text{ that
are continuous functions of time }
\end{eqnarray}%
\begin{eqnarray}
&&1.\int_{t_{u}}^{t_{f}}\left\{ \left\langle \Psi \left( t\right) \left\vert
\left( i\frac{\partial }{\partial t}-H\right) \delta \Psi \left( t\right)
\right. \right\rangle \right. \\
&&\qquad \left. {}\right. \ \ \ \ \ \ \ \ \ \ \ \ \ \ \ \ \ \ \left. -%
\mathcal{L}\left( \Psi \left( t\right) \right) \left\langle \Psi \left(
t\right) \left\vert \delta \Psi \left( t\right) \right. \right\rangle
\right\} dt=0\quad \forall \delta \Psi \left( t\right)
\end{eqnarray}%
\begin{equation}
\left\langle \Phi \left( t\right) |\Phi \left( t\right) \right\rangle =
\end{equation}%
\begin{equation}
\eta \left( t\right)
\end{equation}%
\begin{equation}
\bullet
\end{equation}%
The four centre integral terms can be term can be further developed to give 
\begin{eqnarray}
\sum_{j_{1,}j_{2,}j_{3,}j_{4}=1}^{r}\mathcal{D}%
_{2i_{1}i_{2}j_{1}j_{2}j_{3}j_{4}}\left\langle
j_{1}j_{2}||j_{3}j_{4}\right\rangle ^{a}
&=&\sum_{j_{1,}j_{2,}j_{3,}j_{4}=1}^{r}\sum_{1\leq a,b\leq s}\mathcal{D}%
_{2i_{1}i_{2}j_{1}j_{2}j_{3}j_{4}}\mathcal{R}_{j_{1}j_{2}ab}^{\dagger }%
\mathcal{S}_{abj_{4}j_{3}} \\
&=&\sum_{1\leq a,b\leq s}\mathcal{X}_{i_{1}i_{2}ab}\mathcal{Y}_{ab}
\end{eqnarray}%
where 
\begin{equation}
\mathcal{X}_{i_{1}i_{2}abj_{3}j_{4}}=\sum_{j_{1,}j_{2,}=1}^{r}\mathcal{D}%
_{2i_{1}i_{2}j_{1}j_{2}j_{3}j_{4}}\mathcal{R}_{j_{1}j_{2}ab}^{\dagger }
\end{equation}%
and 
\begin{equation}
\mathcal{Y}_{abj_{3}j_{4}}=\sum_{j_{3,}j_{4}=1}^{r}\mathcal{D}%
_{2i_{1}i_{2}j_{1}j_{2}j_{3}j_{4}}\mathcal{S}_{j_{4}j_{4}ab}
\end{equation}%
In matrix form over the decomposition $\ker L_{N}^{1}\oplus \mathcal{V}_{%
\tilde{\Gamma}}$ with the super vector $\left\vert \tilde{\Gamma}\left( 
\mathfrak{d}_{1}\right) \right) $ being part of the basis of $\mathcal{V}_{%
\tilde{\Gamma}}$ these transformations have the schematic appearence%
\begin{eqnarray}
&&\hat{T}\left( 
\begin{array}{cc}
\left\vert \mathfrak{d}_{ref}\right) ,\left\vert k_{2}\right) ,\ldots
,\left\vert k_{r_{1}}\right) & \left\vert \tilde{\Gamma}\left( \mathfrak{d}%
_{1}\right) \right) \mathfrak{,}\left\vert \mathfrak{x}_{2}\right) ,\ldots
,\left\vert \mathfrak{x}_{r_{2}}\right)%
\end{array}%
\right) =  \notag \\
&&\qquad \qquad \left( 
\begin{array}{cc}
\left\vert \mathfrak{d}_{ref}\right) ,\left\vert k_{2}\right) ,\ldots
,\left\vert k_{r_{1}}\right) & \left\vert \tilde{\Gamma}\left( \mathfrak{d}%
_{1}\right) \right) \mathfrak{,}\left\vert \mathfrak{x}_{2}\right) ,\ldots
,\left\vert \mathfrak{x}_{r_{2}}\right)%
\end{array}%
\right)  \notag \\
&&\times \\
&&\qquad \qquad \qquad \qquad \qquad \qquad \qquad \qquad \left( 
\begin{array}{cc}
\mathbb{T} & \left[ \qquad \ast \qquad \right] \\ 
\mathbb{O} & \left[ 
\begin{array}{cccc}
1 & \mathbb{O} & \mathbf{\cdots } & \mathbb{O} \\ 
\mathbb{O} & \mathbb{O} & \cdots & \mathbb{O} \\ 
\mathbb{O} & \mathbb{O} & \mathbb{O} & \mathbb{O}%
\end{array}%
\right]%
\end{array}%
\right)
\end{eqnarray}%
where $\mathbb{T\in }\left[ \mathcal{R}_{\ker L_{N}^{1}}\left( \mathfrak{A}%
_{\ker }\right) \right] $ and $\ast $ arbitrary submatrices.

with the super vector $\left\vert \tilde{\Gamma}\left( \mathfrak{d}%
_{1}\right) \right) $ being part of the basis of $\mathcal{V}_{\tilde{\Gamma}%
}$ these transformations have the schematic appearence%
\begin{eqnarray}
&&\hat{T}\left( 
\begin{array}{cc}
\left\vert \mathfrak{d}_{ref}\right) ,\left\vert k_{2}\right) ,\ldots
,\left\vert k_{r_{1}}\right) & \left\vert \tilde{\Gamma}\left( \mathfrak{d}%
_{1}\right) \right) \mathfrak{,}\left\vert \mathfrak{x}_{2}\right) ,\ldots
,\left\vert \mathfrak{x}_{r_{2}}\right)%
\end{array}%
\right) =  \notag \\
&&\qquad \qquad \left( 
\begin{array}{cc}
\left\vert \mathfrak{d}_{ref}\right) ,\left\vert k_{2}\right) ,\ldots
,\left\vert k_{r_{1}}\right) & \left\vert \tilde{\Gamma}\left( \mathfrak{d}%
_{1}\right) \right) \mathfrak{,}\left\vert \mathfrak{x}_{2}\right) ,\ldots
,\left\vert \mathfrak{x}_{r_{2}}\right)%
\end{array}%
\right)  \notag \\
&&\times \\
&&\qquad \qquad \qquad \qquad \qquad \qquad \qquad \qquad \left( 
\begin{array}{cc}
\mathbb{T} & \left[ \qquad \ast \qquad \right] \\ 
\mathbb{O} & \left[ 
\begin{array}{cccc}
1 & \mathbb{O} & \mathbf{\cdots } & \mathbb{O} \\ 
\mathbb{O} & \mathbb{O} & \cdots & \mathbb{O} \\ 
\mathbb{O} & \mathbb{O} & \mathbb{O} & \mathbb{O}%
\end{array}%
\right]%
\end{array}%
\right)
\end{eqnarray}%
where $\mathbb{T\in }\left[ \mathcal{R}_{\ker L_{N}^{1}}\left( \mathfrak{A}%
_{\ker }\right) \right] $ and $\ast $ arbitrary submatrices.

Noting that changes in the expectation values of one particle observables
are determined by changes in the one particle reduced density matrix, one
can describe the effects of the radiation field, (approximated as an
effective one particle interaction), on the expectation values of any one
particle observable in terms of retarded two particle propagators.

\bigskip

If the evolution starts at $t=0$ then 
\begin{equation}
D^{1}\mathcal{E}\left( \mathbf{z}_{0},\mathbf{\bar{z}}_{0},0\right) =\frac{1%
}{\mathcal{S}\left( \mathbf{z}_{0},\mathbf{\bar{z}}_{0}\right) }\left( \frac{%
\partial \mathcal{F}\left( \mathbf{z}_{0},\mathbf{\bar{z}}_{0},0\right) }{%
\partial \mathbf{z}},\frac{\partial \mathcal{F}\left( \mathbf{z}_{0},\mathbf{%
\bar{z}}_{0},0\right) }{\partial \mathbf{\bar{z}}}\right)
\end{equation}%
and 
\begin{align}
& D^{2}\mathcal{E}\left( \mathbf{z}_{0},\mathbf{\bar{z}}_{0},0\right) \equiv
\left\{ \frac{1}{\mathcal{S}\left( \mathbf{z}_{0},\mathbf{\bar{z}}%
_{0}\right) }\left\{ \frac{\partial ^{2}\mathcal{H}_{0}\left( \mathbf{z}_{0},%
\mathbf{\bar{z}}_{0}\right) }{\partial \mathbf{x}\partial \mathbf{y}}-%
\mathcal{E}_{0}\left( \mathbf{z}_{0},\mathbf{\bar{z}}_{0}\right) \frac{%
\partial ^{2}\mathcal{S}\left( \mathbf{z}_{0},\mathbf{\bar{z}}_{0}\right) }{%
\partial \mathbf{x}\partial \mathbf{y}}\right. \right.  \notag \\
& \qquad \qquad \qquad \qquad \qquad \qquad \qquad +\frac{\partial ^{2}%
\mathcal{F}\left( \mathbf{z}_{0},\mathbf{\bar{z}}_{0},0\right) }{\partial 
\mathbf{x}\partial \mathbf{y}}-\frac{\partial \mathcal{F}\left( \mathbf{z}%
_{0},\mathbf{\bar{z}}_{0},t\right) }{\partial \mathbf{x}}\frac{\partial 
\mathcal{S}\left( \mathbf{z}_{0},\mathbf{\bar{z}}_{0}\right) }{\partial 
\mathbf{y}}  \notag \\
& \qquad \qquad \qquad \qquad \qquad \left. -\frac{\partial \mathcal{S}%
\left( \mathbf{z}_{0},\mathbf{\bar{z}}_{0}\right) }{\partial \mathbf{x}}%
\frac{\partial \mathcal{F}\left( \mathbf{z}_{0},\mathbf{\bar{z}}%
_{0},0\right) }{\partial \mathbf{y}}-\mathcal{F}\left( \mathbf{z}_{0},%
\mathbf{\bar{z}}_{0},0\right) \frac{\partial ^{2}\mathcal{S}\left( \mathbf{z}%
_{0},\mathbf{\bar{z}}_{0}\right) }{\partial \mathbf{x}\partial \mathbf{y}}%
\right\}
\end{align}%
As%
\begin{equation}
\mathcal{F}\left( \mathbf{\bar{z}},\mathbf{z},t\right) =\frac{\left\langle
\left. \mathbf{z}\left( t\right) \right\vert F\left( t\right) \mathbf{z}%
\left( t\right) \right\rangle }{\left\langle \left. \mathbf{z}\left(
t\right) \right\vert \mathbf{z}\left( t\right) \right\rangle }
\end{equation}%
the short time expansion about the critical point and $t=0$ gives 
\begin{align}
\mathcal{F}\left( \mathbf{z}_{0},\mathbf{\bar{z}}_{0},\delta t\right) &
\approx \mathcal{F}\left( \mathbf{z}_{0},\mathbf{\bar{z}}_{0},0\right) +%
\frac{1}{\mathcal{S}\left( \mathbf{z}_{0},\mathbf{\bar{z}}_{0}\right) }%
\left\{ \left( \frac{\partial \left\langle \left. \mathbf{z}_{0}\right\vert
F\left( 0\right) \mathbf{z}_{0}\right\rangle }{\partial \mathbf{z}}\overset{%
\cdot }{\delta \mathbf{z}}+\frac{\partial \left\langle \left. \mathbf{z}%
_{0}\right\vert F\left( 0\right) \mathbf{z}_{0}\right\rangle }{\partial 
\mathbf{\bar{z}}}\overset{\cdot }{\delta \mathbf{\bar{z}}}\right) \delta
t\right.  \notag \\
& \left. +\,\left\langle \left. \mathbf{z}_{0}\right\vert \dot{F}\left(
0\right) \mathbf{z}_{0}\right\rangle \delta t-\mathcal{F}\left( \mathbf{z}%
_{0},\mathbf{\bar{z}}_{0},0\right) \left( \frac{\partial \mathcal{S}\left( 
\mathbf{z}_{0},\mathbf{\bar{z}}_{0}\right) }{\partial \mathbf{z}}\overset{%
\cdot }{\delta \mathbf{z}}+\frac{\partial \mathcal{S}\left( \mathbf{z}_{0},%
\mathbf{\bar{z}}_{0}\right) }{\partial \mathbf{\bar{z}}}\overset{\cdot }{%
\delta \mathbf{\bar{z}}}\right) \delta t\right\}
\end{align}

and 
\begin{equation}
\left( 
\begin{array}{c}
\delta \mathbf{\bar{z}}\left( t\right) \\ 
\delta \mathbf{z}\left( t\right)%
\end{array}%
\right) =\int_{-\infty }^{+\infty }R\left( t-\tau \right) D^{1}\mathcal{F}%
\left( \mathbf{z}_{0},\mathbf{\bar{z}}_{0},t\right) d\tau
\end{equation}%
where the response function is defined as 
\begin{equation}
\mathcal{R}\left( t\right) =\theta _{+}\left( t\right) \left\{ \left( 
\begin{array}{cc}
\mathbf{0} & \frac{i}{\hslash }\mathbf{C}\left( \mathbf{z}_{0},\mathbf{\bar{z%
}}_{0},0\right) \frac{d}{dt} \\ 
-\frac{i}{\hslash }\mathbf{\bar{C}}\left( \mathbf{z}_{0},\mathbf{\bar{z}}%
_{0},0\right) \frac{d}{dt} & \mathbf{0}%
\end{array}%
\right) -D^{2}\mathcal{E}\left( \mathbf{z}_{0},\mathbf{\bar{z}}_{0},t\right)
\right\} ^{-1}
\end{equation}%
Taking the Laplace transform of $R\left( t\right) $ we obtain 
\begin{equation}
\left( 
\begin{array}{c}
\delta \mathbf{\bar{z}}\left( \xi \right) \\ 
\delta \mathbf{z}\left( \xi \right)%
\end{array}%
\right) =R\left( \xi \right) D^{1}F\left( \mathbf{z}_{0},\mathbf{\bar{z}}%
_{0},\xi \right)
\end{equation}%
where 
\begin{equation}
R\left( \xi \right) =\left\{ \left( 
\begin{array}{cc}
\mathbf{0} & -\xi \mathbf{C}\left( \mathbf{z}_{0},\mathbf{\bar{z}}%
_{0},0\right) \\ 
\bar{\xi}\mathbf{\bar{C}}\left( \mathbf{z}_{0},\mathbf{\bar{z}}_{0},0\right)
& \mathbf{0}%
\end{array}%
\right) -D^{2}\mathcal{E}\left( \mathbf{z}_{0},\mathbf{\bar{z}}_{0},0\right)
\right\} ^{-1}
\end{equation}%
The linear response of the 1-particle reduced density matrix has the Laplace
transform 
\begin{align}
& \delta ^{1}\mathbf{\rho }\left( \mathbf{z}_{0},\mathbf{\bar{z}}_{0}\right)
_{ij}\left( \xi \right) =\left( 
\begin{array}{cc}
\frac{\partial \rho \left( \mathbf{z}_{0},\mathbf{\bar{z}}_{0}\right) _{ij}}{%
\partial \mathbf{\bar{z}}} & \frac{\partial \rho \left( \mathbf{z}_{0},%
\mathbf{\bar{z}}_{0}\right) _{ij}}{\partial \mathbf{z}}%
\end{array}%
\right) \times  \notag \\
& \qquad \qquad \qquad \qquad \left\{ \left( 
\begin{array}{cc}
\mathbf{0} & -\xi \mathbf{C}\left( \mathbf{z}_{0},\mathbf{\bar{z}}%
_{0},t\right) \\ 
\bar{\xi}\mathbf{\bar{C}}\left( \mathbf{z}_{0},\mathbf{\bar{z}}_{0},t\right)
& \mathbf{0}%
\end{array}%
\right) -D^{2}\mathcal{E}\left( \mathbf{z}_{0},\mathbf{\bar{z}}_{0},t\right)
\right\} ^{-1}D^{1}\mathcal{F}\left( \mathbf{z}_{0},\mathbf{\bar{z}}_{0},\xi
\right)
\end{align}%
\begin{equation}
\left\{ \left( 
\begin{array}{cc}
\mathbf{0} & \frac{i}{\hslash }\mathbf{C}\left( \mathbf{z}_{0},\mathbf{\bar{z%
}}_{0},0\right) \frac{d}{dt} \\ 
-\frac{i}{\hslash }\mathbf{\bar{C}}\left( \mathbf{z}_{0},\mathbf{\bar{z}}%
_{0},0\right) \frac{d}{dt} & \mathbf{0}%
\end{array}%
\right) -D^{2}\mathcal{E}_{0}\left( \mathbf{z}_{0},\mathbf{\bar{z}}%
_{0},0\right) \right\} \left( 
\begin{array}{c}
\delta \mathbf{\bar{z}}\left( t\right) \\ 
\delta \mathbf{z}\left( t\right)%
\end{array}%
\right) =D^{1}\mathcal{F}\left( \mathbf{z}_{0},\mathbf{\bar{z}}_{0},0\right)
\end{equation}

As an illustration of the preceding forms of the cosets of $U\left( \mathcal{%
H}^{1}\right) $ we can consider their action on an AGP\ state $\left\vert g^{%
\frac{N}{2}}\right\rangle $ determined by a geminal 
\begin{equation}
\left\vert g\right\rangle =\sum_{1\leq j\leq s}c_{j_{1}}\left\vert \eta
_{j}\eta _{j+s}\right\rangle
\end{equation}%
given by 
\begin{equation}
c_{+1}\left( \mathbf{z}_{+}\right) c_{0}\left( \mathbf{\lambda }\right)
c_{+1}\left( \mathbf{z}_{\sigma +}\right) v_{0}\left( \mathbf{\lambda }_{\xi
}\right) \tprod_{1\leq j\leq \rho }\left\{ u_{+1j}\left( \mathbf{z}%
_{+j}\right) u_{0j}\left( \mathbf{\lambda }_{j}\right) \right\} c_{0}\left( 
\mathbf{\lambda }_{\sigma }\right) \left\vert g^{\frac{N}{2}}\right\rangle
\end{equation}%
$\sigma =$ number of non zero, non equal $\left\{ c_{j}\right\} $'s

\begin{description}
\item[$\bullet $] $\omega _{j}=$ number of equal coefficients in the $j^{th}$
set

\item[$\bullet $] $\rho =$ number of sets of equal $\left\{ c_{j}\right\} $'s

\item[$\bullet $] $\xi =$ number of zero coefficients.
\end{description}

The integral kernels of the operators $h_{1}$ and $h_{2}$ are 
\begin{align}
& h_{1}\left( \mathbf{r}_{1}\mathbf{,\xi }_{1};\mathbf{r}_{1}^{\prime }%
\mathbf{,\xi }_{1}^{\prime }\right) =\left\langle \mathbf{r}_{1}\mathbf{,\xi 
}_{1}\left\vert \left( h_{1}\right) \mathbf{r}_{1}^{\prime }\mathbf{,\xi }%
_{1}^{\prime }\right. \right\rangle   \notag \\
& \qquad \qquad \qquad \qquad \left. =\right. \sum_{_{1\leq l\leq N_{nuc}}}%
\frac{Z_{l}}{\left\Vert \mathbf{R}_{l}-\mathbf{r}_{1}\right\Vert }\delta
\left( \mathbf{r}_{1}-\mathbf{r}_{1}^{\prime }\right) \delta \left( \mathbf{%
\xi }_{1}-\mathbf{\xi }_{1}^{\prime }\right) +\nabla _{\mathbf{r}%
_{1}}^{2}\delta \left( \mathbf{r}_{1}-\mathbf{r}_{1}^{\prime }\right)  
\notag \\
& h_{2}\left( \mathbf{r}_{1}\mathbf{,\xi }_{1},\mathbf{r}_{2}\mathbf{,\xi }%
_{2};\mathbf{r}_{1}^{\prime }\mathbf{,\xi }_{1}^{\prime },\mathbf{r}%
_{2}^{\prime }\mathbf{,\xi }_{2}^{\prime }\right) =\left\langle \mathbf{r}%
_{1}\mathbf{,\xi }_{1},\mathbf{r}_{2}\mathbf{,\xi }_{2}\left. \left(
h_{2}\right) \mathbf{r}_{1}^{\prime }\mathbf{,\xi }_{1}^{\prime },\mathbf{r}%
_{2}^{\prime }\mathbf{,\xi }_{2}^{\prime }\right\vert \right\rangle   \notag
\\
& \qquad \qquad \qquad \qquad \left. =\right. \frac{1}{\left\Vert \mathbf{r}%
_{1}-\mathbf{r}_{2}\right\Vert }\left\{ \delta \left( \mathbf{r}_{1}-\mathbf{%
r}_{1}^{\prime }\right) \delta \left( \mathbf{r}_{2}-\mathbf{r}_{2}^{\prime
}\right) \delta \left( \mathbf{\xi }_{1}-\mathbf{\xi }_{1}^{\prime }\right)
\delta \left( \mathbf{\xi }_{2}-\mathbf{\xi }_{2}^{\prime }\right) \right.  
\notag \\
& \qquad \qquad \qquad \qquad \qquad \qquad \qquad \qquad -\left. \delta
\left( \mathbf{r}_{1}-\mathbf{r}_{2}^{\prime }\right) \delta \left( \mathbf{r%
}_{2}-\mathbf{r}_{1}^{\prime }\right) \delta \left( \mathbf{\xi }_{1}-%
\mathbf{\xi }_{2}^{\prime }\right) \delta \left( \mathbf{\xi }_{2}-\mathbf{%
\xi }_{1}^{\prime }\right) \right\} 
\end{align}%
Noting 
\begin{equation}
L_{2}^{1}\left( \mathbb{A}^{\dagger }\right) =L_{2}^{1}\left( \mathbb{A}%
\right) ^{\dagger }
\end{equation}%
we obtain 
\begin{equation}
F\left( \mathbf{\psi }\right) ^{\dagger }=\frac{1}{2}\sum_{1\leq m,n\leq
2s}\left\{ \frac{1}{\left( 2N-1\right) }\left( L_{2}^{1}\left( \mathbb{D}%
^{2}\left\{ \left( \mathbb{H}_{1}\otimes \mathbb{I}_{1}\right) +\mathbb{H}%
_{2}\right\} \right) \right) _{mn}\right\} \left\vert \psi _{m}\right\rangle
\left\langle \psi _{n}\right\vert ,
\end{equation}%
which shows that eq. $\left( \ref{enderiv}\right) $ can be written as 
\begin{equation}
\frac{\delta \mathcal{E}\left( \mathbf{\psi }\right) }{\delta \psi _{\mu
}^{\ast }}=\left\{ F\left( \mathbf{\psi }\right) -F\left( \mathbf{\psi }%
\right) ^{\dagger }\right\} \left\vert \psi _{\mu }\right\rangle 
\end{equation}%
The energy using the direct product reduced Hamiltonian is given by 
\begin{equation}
\mathcal{E}\left( D^{2},\mathbf{\psi }\right) =\sum_{1\leq i,j,k,l\leq
2s}\left\{ \frac{1}{\left( 2N-1\right) }\left\langle \left. \psi
_{i}\right\vert h_{1}\psi _{k}\right\rangle \delta _{jl}+\left( \left. \psi
_{i}\psi _{j}\right\vert h_{2}\psi _{k}\psi _{l}\right) \right\} D_{klij}^{2}
\end{equation}%
leading to%
\begin{align}
& \left\{ \sum_{1\leq m,n\leq 2s}\left\{ h_{2mjkl}\delta
_{in}+h_{2milk}\delta _{jn}\right\} \left\vert \psi _{m}\right\rangle
\left\langle \psi _{n}\right\vert \right\} \left\vert \psi _{\mu
}\right\rangle  \\
& \left\{ \sum_{1\leq m,n\leq 2s}h_{1mk}\delta _{in}\left\vert \psi
_{m}\right\rangle \left\langle \psi _{n}\right\vert \right\} \left\vert \psi
_{\mu }\right\rangle 
\end{align}%
\begin{align}
& \frac{\delta \mathcal{E}\left( \mathbf{\psi }\right) }{\delta \psi _{\mu
}^{\ast }} \\
& =\sum_{1\leq m,n\leq 2s}\sum_{1\leq i,j,k,l\leq 2s}\left\{ \left\{ \frac{1%
}{\left( 2N-1\right) }h_{1mk}\delta _{in}\delta _{jl}+\left\{
h_{2mjkl}\delta _{in}+h_{2milk}\delta _{jn}\right\} \right\}
D_{klij}^{2}\left\vert \psi _{m}\right\rangle \left\langle \psi
_{n}\right\vert \left\vert \psi _{\mu }\right\rangle \right\}  \\
& =\sum_{1\leq m,n\leq 2s}\left\{ \left\{ L_{2}^{1}\left( \left\{ \frac{1}{%
\left( 2N-1\right) }\mathbb{H}_{1}\otimes \mathbb{I}_{1}+2\mathbb{H}%
_{2}\right\} \mathbb{D}^{2}\right) _{mn}\right\} \left\vert \psi
_{m}\right\rangle \left\langle \psi _{n}\right\vert \right\} \left\vert \psi
_{\mu }\right\rangle 
\end{align}

\begin{align}
& \mathcal{E}\left( \mathbb{D}^{2},\mathbf{\psi }\right) =\func{Tr}\left\{
\left\{ \sum_{1\leq i,k\leq 2s}h_{1ik}a_{i}^{\dagger }a_{k}+\frac{1}{4}%
\sum_{1\leq i,j,k,l\leq 2s}\left\langle ij||kl\right\rangle a_{i}^{\dagger
}a_{j}^{\dagger }a_{l}a_{k}\right\} D^{2N}\right\}  \\
& =\func{Tr}\left\{ \sum_{1\leq i,j,k,l\leq 2s}\frac{1}{4}\left\{ \left\{ 
\frac{1}{\left( 2N-1\right) }\left\{ h_{1ik}\delta _{jl}+h_{1jl}\delta
_{ik}-h_{1il}\delta _{jk}-h_{1jk}\delta _{il}\right\} \right\} +\left\langle
ij||kl\right\rangle \right\} a_{i}^{\dagger }a_{j}^{\dagger
}a_{l}a_{k}D^{2N}\right\}   \notag \\
& =\frac{1}{4}\sum_{1\leq i,j,k,l\leq 2s}\left\{ \frac{1}{\left( 2N-1\right) 
}\left\{ h_{1ik}\delta _{jl}+h_{1jl}\delta _{ik}-h_{1il}\delta
_{jk}-h_{1jk}\delta _{il}\right\} +\left\langle ij||kl\right\rangle \right\} 
\func{Tr}\left\{ a_{i}^{\dagger }a_{j}^{\dagger }a_{l}a_{k}D^{2N}\right\} ,
\end{align}%
where the antisymmetric reduced Hamiltonian matrix elements are 
\begin{subequations}
\begin{align}
K_{ijkl}& =\frac{1}{\left( 2N-1\right) }\left\{ h_{1ik}\delta
_{jl}+h_{1jl}\delta _{ik}-h_{1il}\delta _{jk}-h_{1jk}\delta _{il}\right\}
+\left\langle ij||kl\right\rangle   \notag \\
& =\frac{2}{\left( 2N-1\right) }\left\{ h_{1ik}\delta _{jl}-h_{1il}\delta
_{jk}\right\} +h_{2ijkl}-h_{2ijlk},
\end{align}%
in terms of the matrix elements of the coloumb operator in the direct
product basis 
\end{subequations}
\begin{equation}
h_{2ijkl}=\left( \left. \psi _{i}\psi _{j}\right\vert h_{2}\psi _{k}\psi
_{l}\right) 
\end{equation}%
\begin{equation}
D_{klij}^{2}=\func{Tr}\left\{ a_{i}^{\dagger }a_{j}^{\dagger
}a_{l}a_{k}D^{2N}\right\} .
\end{equation}%
The reduced Hamiltonian matrix $\mathbb{K}\left( \mathbf{\psi }\right) $ can
easily seen to be hermitian as 
\begin{align}
K_{klij}^{\ast }& =\frac{2}{\left( 2N-1\right) }\left\{ h_{1ki}^{\ast
}\delta _{lj}-h_{1li}^{\ast }\delta _{kj}\right\} +h_{2klij}^{\ast
}-h_{2lkij}^{\ast }  \notag \\
& =\frac{2}{\left( 2N-1\right) }\left\{ h_{1ik}\delta _{jl}-h_{1il}\delta
_{jk}\right\} +h_{2ijkl}-h_{2ijlk}=K_{ijkl}.
\end{align}%
We can rearrange terms in the energy expression eq. $\left( \ref{RDO_en}%
\right) $ to obtain 
\begin{align}
\mathcal{E}\left( D^{2},\mathbf{\psi }\right) & =\frac{1}{4}\sum_{1\leq
i,j,k,l\leq 2s}\left\{ \frac{2}{\left( 2N-1\right) }\left\{ h_{1ik}\delta
_{jl}D_{klij}^{2}-h_{1il}\delta _{jk}D_{klij}^{2}\right\}
+h_{2ijkl}D_{klij}^{2}-h_{2ijlk}D_{klij}^{2}\right\}   \notag \\
& =\frac{1}{2}\sum_{1\leq i,j,k,l\leq 2s}\left\{ \frac{2}{\left( 2N-1\right) 
}h_{1ik}\delta _{jl}+h_{2ijkl}\right\} D_{klij}^{2}  \notag \\
& =\frac{1}{2}\sum_{1\leq i,j,k,l\leq 2s}\left\{ \frac{1}{\left( 2N-1\right) 
}\left( h_{1ik}\delta _{jl}+h_{1jl}\delta _{ik}\right) +h_{2ijkl}\right\}
D_{klij}^{2}.
\end{align}%
which shows that 
\begin{subequations}
\label{Rham_1elec}
\begin{align}
\mathcal{E}\left( D^{2},\mathbf{\psi }\right) & =\frac{1}{2}\func{Tr}\left\{
\left\{ \frac{1}{\left( 2N-1\right) }\left( \mathbb{H}_{1}\otimes \mathbb{I}%
_{1}+\mathbb{I}_{1}\otimes \mathbb{H}_{1}\right) +\mathbb{H}_{2}\right\} 
\mathbb{D}^{2}\right\}   \label{Rham_1eleca} \\
& =\frac{1}{2}\func{Tr}\left\{ \left\{ \frac{2}{\left( 2N-1\right) }\mathbb{H%
}_{1}\otimes \mathbb{I}_{1}+\mathbb{H}_{2}\right\} \mathbb{D}^{2}\right\} 
\label{Rham_1elecb}
\end{align}%
This displays that wrt the direct product basis we can introduce the non
antisymmetric Reduced Hamiltonian matrix with an explicitly symmetric one
electron part 
\end{subequations}
\begin{equation}
\mathbb{K}_{d}\left( \mathbf{\psi }\right) =\frac{1}{\left( 2N-1\right) }%
\left( \mathbb{H}_{1}\otimes \mathbb{I}_{1}+\mathbb{I}_{1}\otimes \mathbb{H}%
_{1}\right) +\mathbb{H}_{2},
\end{equation}%
or in a more compact non symmetric form 
\begin{equation}
\mathbb{K}_{d}\left( \mathbf{\psi }\right) =\frac{2}{\left( 2N-1\right) }%
\mathbb{H}_{1}\otimes \mathbb{I}_{1}+\mathbb{H}_{2},
\end{equation}%
where FORDM%
\begin{align}
\mathbb{H}_{1}& \leftrightarrow \left\{ \left. h_{1ik}\right\vert 1\leq
i,k\leq 2s\right\}   \notag \\
\mathbb{H}_{2}& \leftrightarrow \left\{ \left. h_{2ijkl}\right\vert 1\leq
i,j,k,l\leq 2s\right\} 
\end{align}%
and express the energy functional as 
\begin{equation}
\mathcal{E}\left( D^{2},\mathbf{\psi }\right) =\frac{1}{2}\func{Tr}\left\{ 
\mathbb{K}_{d}\left( \mathbf{\psi }\right) \mathbb{D}^{2}\right\} .
\label{RDO_en1}
\end{equation}%
where%
\begin{equation}
\mathbb{K}_{d}\left( \mathbf{\psi }\right) =\frac{2}{\left( 2N-1\right) }%
\mathbb{H}_{1}\otimes \mathbb{I}_{1}+\mathbb{H}_{2}  \label{e2}
\end{equation}%
and 
\begin{equation}
K_{dijkl}=\frac{2}{\left( 2N-1\right) }\left\langle \left. \psi
_{i}\right\vert h_{1}\psi _{k}\right\rangle \delta _{jl}+\left( \left. \psi
_{i}\psi _{j}\right\vert h_{2}\psi _{k}\psi _{l}\right) -\left( \left. \psi
_{i}\psi _{j}\right\vert h_{2}\psi _{l}\psi _{k}\right) .  \label{e3}
\end{equation}%
The partial derivatives $\left\{ \frac{\delta \mathcal{E}\left( \mathbf{\psi 
}\right) }{\delta \psi _{\mu }^{\ast }}\right\} $ are obtained by using eqs. 
$\left( \ref{e1}\right) -\left( \ref{e3}\right) $, $\left( \ref%
{1partfuncder_ex}\right) $ and $\left( \ref{2partfuncder_ex}\right) $
producing 
\begin{align}
& \frac{\delta \mathcal{E}\left( \mathbf{\psi }\right) }{\delta \psi _{\mu
}^{\ast }}  \notag \\
& \left. =\right. \sum_{1\leq m,n\leq 2s}\sum_{1\leq i,j,k,l\leq 2s}\left\{
\left\{ \frac{1}{2N-1}h_{1mk}\delta _{in}\delta _{jl}+h_{2mjkl}\delta
_{in}\right\} \overset{}{D_{klij}^{2}\left\vert \psi _{m}\right\rangle
\left\langle \psi _{n}\right\vert }\right\} \left\vert \psi _{\mu
}\right\rangle .
\end{align}%
This allows us to define the generalized Fock operator
\begin{subequations}
\begin{align}
F\left( \mathbf{\psi }\right) & =\sum_{1\leq m,n\leq 2s}\left\{ \sum_{1\leq
i,j,k,l\leq 2s}\left\{ \frac{1}{\left( 2N-1\right) }h_{1mk}\delta
_{in}\delta _{jl}+h_{2mjkl}\delta _{in}\right\} D_{klij}^{2}\right\}
\left\vert \psi _{m}\right\rangle \left\langle \psi _{n}\right\vert  \\
& =\sum_{1\leq m,n\leq 2s}\left\{ \left( L_{2}^{1}\left( \left\{ \frac{1}{%
\left( 2N-1\right) }\mathbb{H}_{1}\otimes \mathbb{I}_{1}+\mathbb{H}%
_{2}\right\} \mathbb{D}^{2}\right) \right) _{mn}\right\} \left\vert \psi
_{m}\right\rangle \left\langle \psi _{n}\right\vert .
\end{align}%
\end{subequations}
\begin{equation}
\mathbb{F}_{2}\left( \mathbf{\psi }\right) =\frac{1}{\left( 2N-1\right) }%
\mathbb{H}_{1}\otimes \mathbb{I}_{1}+\mathbb{H}_{2}
\end{equation}%
\begin{equation}
\mathbb{F}\left( \mathbf{\psi }\right) =L_{2}^{1}\left( \mathbb{F}_{2}\left( 
\mathbf{\psi }\right) \mathbb{D}^{2}\right) 
\end{equation}%
In case (b) the energy is given by 

\begin{equation}
\mathcal{E}\left( D^{2},\mathbf{\psi }\right) =\func{Tr}\left\{ \mathbb{H}%
_{1}\mathbb{D}^{1}\right\} +\frac{1}{4}\func{Tr}\left\{ \mathbb{V}_{2}%
\mathbb{D}^{2}\right\} 
\end{equation}

\end{document}
